\chapter{Grover's Algorithm and Amplitude Amplification}
\label{ch:lec03}

\begin{objectives}
\begin{itemize}[leftmargin=1.4em,nosep]
  \item Trace the historical arc from Deutsch's algorithm to unstructured search.
  \item Derive the Grover iteration geometrically.
  \item Generalize to amplitude amplification and state its quadratic speedup.
  \item Sketch the query lower bound for unstructured search.
\end{itemize}
\end{objectives}

\todo{Write the introductory paragraph orienting the lecture.}

\section{Deutsch's algorithm}
\label{sec:lec03-1}
\todo{Draft this section.}

\section{The unstructured search problem}
\label{sec:lec03-2}
\todo{Draft this section.}

\section{The Grover iteration}
\label{sec:lec03-3}
\todo{Draft this section.}

\section{Amplitude amplification}
\label{sec:lec03-4}
\todo{Draft this section.}

\section{Query complexity lower bounds}
\label{sec:lec03-5}
\todo{Draft this section.}

\exercisesection
\begin{exercise}
\todo{Add exercise.}
\end{exercise}
\begin{exercise}
\todo{Add exercise.}
\end{exercise}
\begin{exercise}
\todo{Add exercise.}
\end{exercise}

\readingsection
\begin{itemize}[leftmargin=1.4em]
  \item Lin Lin, \emph{qasc notes}, Ch.~2.
  \item Lin Lin, \emph{live notes}, Ch.~11 (amplitude amplification).
  \item A.~Childs, \emph{Lecture Notes on Quantum Algorithms} (Grover, lower bounds).
\end{itemize}
